\paragraph{В.2. Вероятность}

\subparagraph{B.2.1}
Профессор Ванстоун бросает симметричную монету один раз, а профессор Уноппетри -- два раза. 
Чему равна вероятность того, что профессор Ванстоун получит больше выпавших орлов,
чем профессор Уноппетри?

\textbf{Ответ:}

Вероятность того что профессору Уноппетри выпадет два решки $1/4$.
Вероятность того что профессору Ванстоуну выпадет орёл $1/2$.
События $A$ и $B$ независимые - $\Pr \{ A \cap B \} = \Pr \{ A \} \Pr\{B \} = 1/8$


\subparagraph{B.2.2}
Докажите \textbf{неравенство Буля}: для любой конечной или бесконечной счетной 
последовательности событий \( A_1, A_2, \dots \) справедливо следующее неравенство:
\[
\Pr \{ A_1 \cup A_2 \cup \dots \} \leq \Pr \{ A_1 \} + \Pr \{ A_2 \} + \dots .
\]

\textbf{Ответ:}

\subparagraph{B.2.3}
Имеется тщательно перетасованная колода из десяти карт, пронумерованных числами от 1 до 10. 
Из колоды последовательно вынимаются три карты. 
Какова вероятность того, что эти карты будут находиться в порядке возрастания их достоинства?

\subparagraph{B.2.4}
Докажите, что
\[
\Pr \{ A \mid B \} + \Pr \{ \overline{A} \mid B \} = 1.
\]

\subparagraph{B.2.5}
Докажите, что для любого набора событий \( A_1, A_2, \dots, A_n \)
\[
\Pr \{ A_1 \cap A_2 \cap \dots \cap A_n \} = \Pr \{ A_1 \} 
\cdot \Pr \{ A_2 \mid A_1 \} \cdot \Pr \{ A_3 \mid A_1 \cap A_2 \} \cdots 
\Pr \{ A_n \mid A_1 \cap A_2 \cap \dots \cap A_{n-1} \}.
\]